\section{Region of interest}
\label{sec:roi}

The region of interest tells RAFTcorr which points to measure. It is a mask on frame~1, and
every result is reported on its points: the result arrays cover its bounding box, and the points
outside the region stay empty. A run does not start until frame~1 has one. The network is given
the bounding box of the region plus a margin (Section~\ref{sec:parameters}); a region drawn
tightly around the specimen keeps that box, and the memory each pair needs, small.

Any other frame can hold a mask of its own. Such a mask is drawn on the image of that frame,
where the specimen is at that moment, and it can leave out a crack or a hole there or serve as
the region of interest of a reference frame (Section~\ref{sec:roi-later-frames}).

\subsection{The hint}

The box at the top of the \ui{REGION OF INTEREST} section says which frames need a region of
interest. It follows the tracking mode (Section~\ref{sec:workflow}):
\begin{itemize}
  \item Before any images are loaded: \ui{Load images first, then draw a Region of Interest on
        frame 1.}
  \item \ui{Accumulative mode}: only frame~1 needs a region of interest. \ui{All later frames
        are compared against it directly.}
  \item \ui{Incremental}: only frame~1 needs a region of interest. RAFTcorr carries it to the
        reference frames, which the hint lists (the first eight, then an ellipsis), with the
        measured displacement. \ui{Draw on another frame only to replace that region or to
        leave out cracks and holes.}
\end{itemize}

\subsection{The toolbar}

\screenshot[0.45\textwidth]{fig-roi-toolbar.png}{The \ui{REGION OF INTEREST} section with the
hint, the toolbar and the menu of \ui{+ Add} open: \ui{Polygon}, \ui{Rectangle}, \ui{Circle} and
\ui{Circle (3-point)}.}{fig:roi-toolbar}

Below the hint are three rows of buttons (Figure~\ref{fig:roi-toolbar}): \ui{+ Add} and
\ui{Cut} (the button with the scissors); \ui{Import} and \ui{Batch Import}; \ui{Save},
\ui{Invert} and \ui{Clear}. Except \ui{Batch Import}, every button acts on the mask of the frame
shown on the canvas, and drawing, importing or clearing opens that frame for editing
(Section~\ref{sec:roi-editing}).

\subsubsection*{Drawing shapes}

\ui{+ Add} and \ui{Cut} open the same menu of four shapes. A shape picked from \ui{+ Add} is
added to the mask; a shape picked from \ui{Cut} is removed from it, which is how you leave out
holes, voids and damaged areas. Then draw the shape on the canvas:
\begin{description}[leftmargin=1.2cm,style=nextline]
  \item[Polygon]
    Click each vertex in turn; double-click to place the last vertex and close the polygon, or
    press \kbd{Enter} to close it after the last vertex you clicked.
    A polygon needs at least three vertices; with fewer, nothing is drawn.
  \item[Rectangle]
    Press at one corner, drag to the opposite corner and release.
  \item[Circle]
    Press at the centre, drag outward to set the radius and release.
  \item[Circle (3-point)]
    Click three points on the edge of the circle. After the second click a preview circle runs
    through the two points and the cursor; the third click fixes the circle that passes through
    all three. Use it when the edge of a round feature is visible but its centre is not. If the
    three points lie almost on one line, no circle is drawn and the log asks for points spread
    around the edge.
\end{description}

While you draw, the cursor is a cross, the outline is blue for \ui{+ Add} and red for \ui{Cut},
and the button you started from is highlighted. A shape is applied as soon as it is complete.
Each pick from the menu draws one shape; pick again for the next one. If you pick the wrong
shape, pick the right one; it replaces the first. To abandon a shape, press \kbd{Esc} before it
is complete, whether you have started drawing it or only picked it from the menu. Shapes that reach beyond the image are cut at its border.

\begin{notebox}
The toolbar has no brush and no undo. Before a change you may want to take back, click
\ui{Save} to keep a copy of the mask, and \ui{Import} it to return to it.
\end{notebox}

\subsubsection*{Import, Save, Invert and Clear}

\begin{description}[leftmargin=1.2cm,style=nextline]
  \item[Import]
    Opens the \ui{Import Mask Image} dialog (PNG, BMP, TIFF or JPEG). The file replaces the mask
    of the frame shown. Bright pixels are inside: a pixel counts as inside when it is brighter
    than mid-grey (127 on a scale of 0 to 255), so masks stored as 0 and 255 and masks stored as
    0 and 1 both work. A colour file is read as its luminance. A mask of another size than the
    images is resized to the image size (nearest neighbour). If the file cannot be read, the
    mask stays as it was, and the log says \emph{``Could not read the mask file''} with the file
    name and the reason.
  \item[Batch Import]
    Assigns mask files to many frames at once (Section~\ref{sec:roi-batch}).
  \item[Save]
    Writes the mask of the frame shown to a PNG file, white (255) inside and black (0) outside.
    The dialog proposes a name made of the image folder and the frame number, next to the image
    folder rather than inside it: for frame~1 of the folder \file{specimen}, the file
    \file{specimen\_roi\_frame1.png} beside \file{specimen}. A mask saved among the images would be
    loaded as an image the next time you open the folder. \ui{Import}, \ui{Batch Import}
    and \file{raftcorr run --roi} (Section~\ref{sec:cli}) read this format.
  \item[Invert]
    Swaps inside and outside of the mask of the frame shown, whether you drew it or it came from
    \ui{Import}, \ui{Batch Import}, the image list or a project file. On frame~1 without a mask it
    selects the whole image. On a later frame without a mask of its own, while frame~1 has one,
    nothing changes, and the log explains that the region of frame~1 is used there.
  \item[Clear]
    Removes the mask of the frame shown. The masks of the other frames stay; to clear several
    frames at once, select them in the image list and right-click (Section~\ref{sec:images}).
\end{description}

\subsection{Editing mode}
\label{sec:roi-editing}

\screenshot[0.5\textwidth]{fig-roi-canvas.png}{The canvas in editing mode for frame~1: the
banner \ui{EDITING REGION OF INTEREST FOR FRAME 01} above the image, and the region of interest
in blue over the gauge section of a plate with two holes, drawn as a rectangle with two circles
cut out.}{fig:roi-canvas}

While you edit a mask, a yellow banner (amber in the light theme) above the canvas names the
frame: \ui{EDITING REGION OF INTEREST FOR FRAME 01} (Figure~\ref{fig:roi-canvas}). The canvas
shows the image of that frame with its mask as a blue tint.

You enter editing mode with the \ui{Region} button of a frame in the image list (\ui{Need},
\ui{Edit} or \ui{Add}), which also shows that frame, or by drawing, importing or clearing with
the toolbar, which edits the frame shown. To edit another frame, click its \ui{Region} button, or move to it with
the image list or the frame slider: the banner follows. Every change is stored at once; there is
nothing to confirm. Editing mode ends when a run starts. Once results exist, moving to another
frame also ends it, and the canvas shows the displacement field again.

A typical region of interest for a plate with two holes (Figure~\ref{fig:roi-canvas}):
\begin{enumerate}
  \item Choose the tracking mode first (Section~\ref{sec:workflow}).
  \item Show frame~1 and choose \menu{+ Add}{Rectangle}. Drag a rectangle over the gauge
        section.
  \item Choose \menu{Cut}{Circle} and drag a circle over each hole, or use
        \menu{Cut}{Circle (3-point)} and click three points on the edge of each hole.
  \item Check the blue region. Click \ui{Save} to keep a copy for later runs or for the command
        line.
\end{enumerate}

\subsection{Masks on later frames}
\label{sec:roi-later-frames}

A mask on a frame other than frame~1 is drawn, imported or batch-imported on the image of that
frame, in its own coordinates. RAFTcorr reads the mask of every frame that has one
(Table~\ref{tab:frame-masks}). In every tracking mode, a point of the region of interest whose
position in frame $k$ falls outside the mask of frame $k$ is left empty in frame $k$: this is the
way to remove a crack, a hole that opens, or a part of the image where the speckle is lost. In
incremental tracking, the mask of a reference frame also replaces the region of frame~1 carried
to that frame, and the pairs compared with it are solved over the mask.

\begin{table}[htbp]
  \centering
  \caption{What the mask of frame $k$ (not frame~1) does.}
  \label{tab:frame-masks}
  \begin{tabularx}{\textwidth}{>{\raggedright\arraybackslash}p{3.2cm}
                               >{\raggedright\arraybackslash}X
                               >{\raggedright\arraybackslash}X}
    \toprule
    Tracking Mode & Frame $k$ has a mask & Frame $k$ has no mask \\
    \midrule
    \ui{Accumulative}
      & Points that fall outside it are empty in frame $k$ only.
      & Nothing is left out. \\
    \ui{Incremental}, frame $k$ is a reference frame
      & Points outside it are empty in frame $k$ and in every frame measured from it; the pairs
        compared with frame $k$ are solved over the mask.
      & The region of frame~1 is carried to frame $k$ with the displacement measured so far. \\
    \ui{Incremental}, frame $k$ is not a reference frame
      & Points that fall outside it are empty in frame $k$.
      & Nothing is left out. \\
    \bottomrule
  \end{tabularx}
\end{table}

Three rules follow from this:
\begin{itemize}
  \item \textbf{A mask on a later frame covers everything that stays.} It is not a mask of the
        crack; it is the whole specimen as it appears in that frame, minus the crack. On a frame
        without a mask of its own, drawing starts from an empty mask, so a \ui{Cut} alone leaves
        nothing. Add a shape over the specimen first, then cut the crack out, or import masks
        made by a segmentation tool with \ui{Batch Import}.
  \item \textbf{A mask cannot add points.} Results exist only for the points of the region of
        interest of frame~1. A mask on a later frame that is larger than the specimen adds
        nothing.
  \item \textbf{In accumulative tracking, a mask acts on its own frame only.} A crack that stays
        open needs a mask on every frame where it is open. \ui{Batch Import} assigns one mask
        file to many frames at once.
\end{itemize}

\begin{tipbox}
In incremental tracking the region of frame~1 follows the specimen from one reference frame to
the next. A specimen that moves or rotates across the image does not need masks on the
reference frames; give one only where you want to leave something out.
\end{tipbox}

\subsection{Batch Import}
\label{sec:roi-batch}

\screenshot[\textwidth]{fig-batch-import.png}{The \ui{Batch Import Region of Interest Masks}
dialog: the \ui{Available Masks} of the chosen folder (left), the \ui{Frame Assignments} with
their \ui{Frame}, \ui{Image} and \ui{Mask} columns (middle), and the \ui{Preview} of the selected
frame with its mask (right).}{fig:batch-import}

\ui{Batch Import} opens the \ui{Batch Import Region of Interest Masks} dialog
(Figure~\ref{fig:batch-import}). It lists every frame, because RAFTcorr reads the mask of every
frame.
\begin{description}[leftmargin=1.2cm,style=nextline]
  \item[Mask Folder]
    \ui{Browse...} chooses the folder of the mask files. Its image files (PNG, BMP, TIFF, JPEG,
    JPEG~2000 and WebP) appear under \ui{Available Masks}, in natural order (\file{mask2.png}
    before \file{mask10.png}). A mask whose size differs from the
    images is greyed out and cannot be assigned; its tooltip gives both sizes, and a line above
    the lists counts these masks. Unlike \ui{Import}, the dialog does not resize masks.
  \item[Auto-Match by Name]
    Gives each image the mask whose file name ends in the same number: \file{mask6.png} goes to
    \file{image6.tif}, whatever frame that image is. If several images end in that number, each
    of them gets the mask. Only when the image names carry no number is the last number of the
    mask name read as the frame position counted from 0: \file{mask0.png} goes to frame~1.
  \item[Assign Sequential]
    Gives the masks, in the order of the list, to frames 1, 2, 3, and so on.
  \item[Assign Selected ->]
    Pairs the masks selected on the left with the frames selected in the middle: one mask with
    one frame, one mask with many frames, or as many masks as frames in the order of selection.
    Several masks for one frame are refused.
  \item[Clear All]
    Removes every assignment. \ui{Auto-Match by Name} and \ui{Assign Sequential} also start
    from no assignments; \ui{Assign Selected ->} adds to the existing ones.
  \item[Preview]
    Shows the image of the frame selected in the middle with its assigned mask, or the mask
    selected on the left over that image. \ui{View:} switches between \ui{Image only},
    \ui{Image + Mask} (default) and \ui{Mask only}; \ui{Alpha:} sets the opacity of the mask
    on a slider from 0 to 100 (default 50); \ui{Mask color:} offers \ui{Blue}, \ui{Red},
    \ui{Green} and \ui{Yellow}. The mouse wheel zooms and dragging pans.
\end{description}
\ui{OK} reads the assigned files with the same rule as \ui{Import} (bright is inside) and
replaces the masks of those frames; the log lists every frame that received one. \ui{Cancel}
changes nothing.

\begin{warnbox}
When the image names carry no number, \ui{Auto-Match by Name} counts frames from 0, as data
files do, and masks numbered from 1 land one frame late. Use \ui{Assign Sequential} or
\ui{Assign Selected ->} instead, and check the \ui{Mask} column before you click \ui{OK}.
\end{warnbox}

A mask with no white pixel, assigned to frame~1, leaves frame~1 without a region of interest:
the run then does not start (Section~\ref{sec:running}).
